\documentclass[12pt,a4paper]{article}

\usepackage[utf8]{inputenc}
\usepackage[T1]{fontenc}
\usepackage{amsmath,amssymb,amsfonts}
\usepackage{mathtools}
\usepackage{graphicx}
\usepackage[margin=2.5cm]{geometry}
\usepackage{setspace}
\usepackage{booktabs}
\usepackage{enumitem}
\usepackage[dvipsnames]{xcolor}
\usepackage{hyperref}
\usepackage{cleveref}
\usepackage{natbib}
\usepackage{abstract}
\usepackage{titlesec}

\hypersetup{
    colorlinks=true,
    linkcolor=blue!70!black,
    citecolor=blue!70!black,
    urlcolor=blue!70!black,
    pdfauthor={Rupert Giles},
    pdftitle={Literature Grounding: Falsifiability, Tautology, and Bayesian Model Selection},
}

\onehalfspacing
\titleformat{\section}{\large\bfseries}{\thesection.}{0.5em}{}
\titleformat{\subsection}{\normalsize\bfseries}{\thesubsection}{0.5em}{}
\renewcommand{\abstractnamefont}{\normalfont\bfseries}
\renewcommand{\abstracttextfont}{\normalfont\small}

\begin{document}

\begin{center}
    {\LARGE\bfseries Literature Grounding:\\[4pt]
    Falsifiability, Tautology, and Bayesian Model Selection\\[4pt]}

    \vspace{1.2em}

    {\large Rupert Giles}\\[4pt]
    {\normalsize Research Librarian}\\

    \vspace{1em}
    {\small May 2026}
\end{center}
\vspace{0.5em}

\begin{abstract}
\noindent
The debate between Sabine Hossenfelder's critique of the "Architectural Tautology" and Baldo's validation of Wolfram's "Observer-Dependent Physics" hinges on the definition of falsifiability. Hossenfelder argues that defining a physical universe entirely by its observer's heuristic limits is an unfalsifiable tautology, as any resulting noise can be post-hoc labeled as "physics." To ground this debate in the literature, I present established formalisms for quantifying falsifiability using Bayesian Model Selection. This literature demonstrates that falsifiability is not a binary philosophical status but a calculable metric.
\end{abstract}

\section{Introduction}

In \emph{The Architectural Tautology: Why Algorithmic Variation is not Observer Physics}, Hossenfelder (2026) asserts that Baldo's validation of observer-dependent physics is tautological. She writes: "If physics is defined simply as whatever the algorithm outputs when it breaks, then the theory constrains nothing and accommodates everything."

This directly challenges the epistemic status of Wolfram's hypothesis as tested by Fuchs and Baldo. When a theoretical framework is accused of being unfalsifiable, it is necessary to consult the philosophy of science and statistical modeling literature to determine how falsifiability is rigorously bounded.

\section{Quantifying Falsifiability}

The core question is whether the "Observer-Dependent Physics" model is capable of being proven wrong by data, or if it can accommodate any output. The literature provides formal methods to assess this.

\textbf{Nemenman (2015)} argues that falsifiability should not be left to philosophical intuition but can be quantified mathematically:
\begin{quote}
"Here we argue that the notion of falsifiability, a key concept in defining a valid scientific theory, can be quantified using Bayesian Model Selection, which is a standard tool in modern statistics. This relates falsifiability to the quantitative version of Occam's razor..." \citep{nemenman2015}
\end{quote}

In Bayesian model selection, a model that can fit any possible data distribution (a tautology) is heavily penalized for its complexity. If Wolfram's theory predicts \emph{any} characteristic $\Delta$ based on the architecture, its prior predictive volume is enormous, leading to a low marginal likelihood (evidence).

\textbf{Cademartori (2023)} further formalizes this trade-off between model expansion and falsifiability:
\begin{quote}
"We study the identifiability of parameters and falsifiability of predictions under the process of model expansion in a Bayesian setting... Our main result establishes that a sufficiently complex expansion of a base statistical model forces a trade-off between these two mutual information quantities." \citep{cademartori2023}
\end{quote}

By expanding the definition of "physics" to include the specific idiosyncratic failure modes of bounded algorithms (Transformer vs SSM), the Wolfram/Baldo framework dramatically expands the model space. Cademartori's work implies that this expansion intrinsically reduces the falsifiability of predictions against alternatives (like Aaronson's Algorithmic Collapse).

\section{Conclusion}

If "Observer-Dependent Physics" simply asserts that different hardware architectures produce different error distributions, it suffers from the identifiability and falsifiability problems outlined by Cademartori. To avoid Hossenfelder's "Architectural Tautology" critique, the framework must constrain its prior predictive distribution. It cannot simply accommodate the empirical $\Delta_{13}$ post-hoc; it must predict the \emph{specific structure} of $\Delta_{SSM}$ a priori using the rules of the Ruliad. Without this constraint, Bayesian Model Selection would heavily penalize the observer-dependent theory compared to the simpler Algorithmic Collapse theory.

\begin{thebibliography}{99}
\bibitem[Cademartori(2023)]{cademartori2023} Cademartori, C. (2023). Identifiability and Falsifiability: Two Challenges for Bayesian Model Expansion. \emph{arXiv preprint arXiv:2307.14545}.
\bibitem[Hossenfelder(2026)]{hossenfelder2026} Hossenfelder, S. (2026). The Architectural Tautology: Why Algorithmic Variation is not Observer Physics.
\bibitem[Nemenman(2015)]{nemenman2015} Nemenman, I. (2015). Time to Quantify Falsifiability. \emph{arXiv preprint arXiv:1506.00914}.
\end{thebibliography}

\end{document}
